\renewcommand{\Anno}{2022}

\renewcommand{\NumeroProblema}{1}

\renewcommand{\Testo}{
Nel 1997 gli allievi del Settore di Ingegneria della Scuola Samt'Anna si recano in viaggio di istruzione a New York City. Thomas, il più estroso del gruppo, ispirandosi alla geometria ed alla viabilità locali, introduce il concetto di $n$-Manhattan, cioè di una città che si può schematizzare come un reticolo costituito da $2n$ strade, di cui $n$ in direzione Nord-Sud ed $n$ in direzione Est-Ovest, che si incontrano formando $n^{2}$ incroci. Tutte le strade sono a senso unico in maniera alternata. Ad esempio, nella figura seguente è rappresentata una 5-Manhattan.

\begin{center}

\psset{unit=0.52ex}
\pspicture(40,42)

\multido{\nA=0+10}{5}{\multido{\nB=0+10}{5}{\psdots[dotscale=1.5](\nA,\nB)}}
\multido{\n=0+10}{5}{\psline(\n,0)(\n,40)\psline(0,\n)(40,\n)}
\multido{\nA=0+20}{3}{\multido{\nB=0+10}{4}{
\psline[arrowscale=2]{->}(\nB,\nA)(!\nB \space 6.5 add \nA)
\psline[arrowscale=2]{->}(\nA,\nB)(!\nA\space \nB \space 6.5 add)}
}
\multido{\nA=10+20}{2}{\multido{\nB=40+-10}{4}{
\psline[arrowscale=2]{->}(\nB,\nA)(!\nB \space 6.5 sub \nA)
\psline[arrowscale=2]{->}(\nA,\nB)(!\nA\space \nB \space 6.5 sub)}
}
\rput[tr](-1.5,-1.5){$A$}
\rput[br](-1.5,41.5){$B$}
\rput[bl](12,12){$C$}

\endpspicture

\end{center}

Diciamo che il \emph{costo} per andare da un incrocio $X$ ad un incrocio $Y$ è $k$ se il percorso più corto che va da $X$ a $Y$ rispettando i sensi unici attraversa $k$ tratti del reticolo. Ad esempio, nella figura il costo per andare da $B$ a $C$ è 4, e anche il costo per andare da $A$ a $C$ è 4.

\begin{enumerate}
\renewcommand{\labelenumi}{(\alph{enumi})}

\item  Siano $X$ e $Y$ due incroci in una 1000-Manhattan, lontani dal perimetro esterno. Sappiamo che il costo per andare da $X$ a $Y$ è 10.

Determinare quanto vale al massimo il costo per andare da $Y$ a $X$.

\item  Consideriamo, in una 1000-Manhattan, l'incrocio $A$ posto in un vertice del perimetro esterno da cui partono due strade (come nel vertice in basso a sinistra nella figura).

Determinare quanti sono gli incroci raggiungibili a partire da $A$ con un costo minore o uguale di 999.

\item  Consideriamo, in una 1000-Manhattan, un incrocio $B$ posto in un vertice del perimetro esterno da cui parte esattamente una strada (come nel vertice in alto a sinistra nella figura).

Determinare quanti sono gli incroci raggiungibili a partire da $B$ con un costo minore o uguale di 999.

\end{enumerate}
}

\renewcommand{\Risposta}{}

\renewcommand{\Aiutino}{}

\renewcommand{\SoluzioneUfficiale}{}

\renewcommand{\SoluzioniAlternative}{}

\renewcommand{\Osservazioni}{}

\renewcommand{\Commenti}{}
